\subsection{Derivations on Semisimple Algebras}

In this subsection and the next ones, K is a commutative ring, A a K-algebra, B the K-algebra \(A \otimes_{K} A^{o}\) , and \(\varepsilon\) the K-linear mapping from B to A defined by \(\varepsilon(x \otimes y) = xy\) for x, y in A.

Recall (III, §4, No.~3, p.~467) that every (A, A)-bimodule P can be viewed as a left B-module, with external law given by \((a \otimes a')p = apa'\) for \(a, a'\) in A and p in P. Conversely, every B-module can be viewed as an (A, A)-bimodule. We endow A with its canonical (A, A)-bimodule structure and with the corresponding B-module structure; we endow B with the (A, A)-bimodule structure corresponding to the B-module \(B_s\) . We therefore have

\[
a (x \otimes y) a ^ {\prime} = (a \otimes a ^ {\prime}) (x \otimes y) = a x \otimes y a ^ {\prime}
\]

for \(a, a', x, y\) in A, where the product \(ya'\) is calculated in the algebra A.

The K-linear mapping \(\varepsilon\) is a homomorphism of (A, A)-bimodules.

PROPOSITION 5. --- The following properties are equivalent:

\begin{enumerate}
\def\labelenumi{(\roman{enumi})}
\item
  The B-module A is projective.
\item
  There exists an element \(e\) of the (A,A)-bimodule B satisfying the following two conditions: \(\varepsilon(e) = 1\) and \(ae = ea\) for every \(a \in \mathrm{A}\) .
\end{enumerate}

The mapping \(\varepsilon: \mathrm{B} \to \mathrm{A}\) is surjective because we have \(\varepsilon(a \otimes 1) = a\) for every \(a \in \mathrm{A}\) ; it is a homomorphism of (A, A)-modules, hence a B-linear mapping. If the B-module A is projective, then there exists a section \(s\) of \(\varepsilon\) (II, §2, No.~2, p.~231, Proposition 4); it is a homomorphism of (A, A)-bimodules from A to B. If we set \(e = s(1)\) , then we have \(\varepsilon(e) = \varepsilon(s(1)) = 1\) and \(ae = s(a) = ea\) for every \(a \in \mathrm{A}\) . So (i) implies (ii).

Conversely, let e be an element of B satisfying the conditions in item (ii). Define a mapping s from A to B by the formula

\[
s (a) = a e = e a.\tag{1}
\]

It is a homomorphism of (A, A)-modules, and we have \(\varepsilon \circ s = 1_{\mathrm{A}}\) ; in other words, \(s\) is a B-linear section of the surjective mapping \(\varepsilon\) . Consequently, the B-module A is isomorphic to the direct factor submodule \(s(\mathrm{A})\) of \(\mathrm{B}_s\) (II, §1, No.~9, p.~211, Proposition 15) and is therefore projective (II, §2, No.~2, p.~231, Proposition 4). This proves that (ii) implies (i).

Remarks. --- 1) Let \(e = \sum_{i=1}^{r} a_i \otimes a_i'\) be an element of B. The conditions in item (ii) of Proposition 5 translate to the formulas

\[
\sum_ {i = 1} ^ {r} a _ {i} a _ {i} ^ {\prime} = 1,\tag{2}
\]

\[
\sum_ {i = 1} ^ {r} a a _ {i} \otimes a _ {i} ^ {\prime} = \sum_ {i = 1} ^ {r} a _ {i} \otimes a _ {i} ^ {\prime} a \quad \text { for   every } a \in A.\tag{3}
\]

When they are satisfied, e is an idempotent in B. Indeed, we then have the relations

\[
e ^ {2} = \sum_ {i = 1} ^ {r} a _ {i} e a _ {i} ^ {\prime} = \sum_ {i = 1} ^ {r} e a _ {i} a _ {i} ^ {\prime} = e.
\]

\begin{enumerate}
\def\labelenumi{\arabic{enumi})}
\setcounter{enumi}{1}
\tightlist
\item
  Let K be a commutative field, let A be a K-algebra, and let M be an A-module. The group \(\operatorname{End}_{\mathrm{K}}(\mathrm{M})\) is endowed with an (A, A)-bimodule structure defined by
\end{enumerate}

\[
a u a ^ {\prime} (x) = a u \left(a ^ {\prime} x\right)
\]

for every \(a, a' \in A\) , every \(u \in \operatorname{End}_{\mathrm{K}}(\mathrm{M})\) , and every \(x \in M\) . We endow it with the associated B-module structure. Let \(e = \sum_{i=1}^{r} a_i \otimes a'_i\) be an element of B satisfying the conditions in part (ii) of Proposition 5, so also relations (2) and (3). If \(p \in \operatorname{End}_{\mathrm{K}}(\mathrm{M})\) is a projector whose image N is an A-submodule of M, then ep is an A-linear projector with the same image.

Indeed, the image of \(ep\) is contained in N. If \(x\) belongs to N, then so does \(a_i' x\) , so that we have \(p(a_i' x) = a_i' x\) and

\[
e p (x) = \sum_ {i = 1} ^ {r} a _ {i} a _ {i} ^ {\prime} x = x
\]

by formula (2). We deduce from formula (3) that \(aep(x) = ep(ax)\) for every \(a \in \mathrm{A}\) and every \(x \in \mathrm{M}\) , which proves that \(ep\) is A-linear.

THEOREM 2. --- Let K be a commutative field and A a K-algebra. The following properties are equivalent:

\begin{enumerate}
\def\labelenumi{(\roman{enumi})}
\item
  The K-algebra A is absolutely semisimple.
\item
  The K-algebra \(\mathrm{B} = \mathrm{A}\otimes_{\mathrm{K}}\mathrm{A}^{\circ}\) is semisimple.
\item
  The B-module A is projective.
\item
  There exists an element \(e\) of the (A,A)-bimodule B satisfying \(\varepsilon(e) = 1\) and \(ae = ea\) for every \(a \in \mathrm{A}\) .
\end{enumerate}

Suppose that the algebra A is absolutely semisimple, hence semisimple. Then the algebra \(A^{o}\) is semisimple (VIII, p.~137, Proposition 2), and it follows from Corollary 1 of VIII, p.~234 that \(B = A \otimes_{K} A^{o}\) is a semisimple K-algebra. Therefore, (i) implies (ii).

Since every module over a semisimple ring is projective (VIII, p.~138, Proposition 4), (ii) implies (iii).

The equivalence of (iii) and (iv) follows from Proposition 5. To complete the proof, let us show that (iv) implies (i). Let \(e = \sum_{i=1}^{r} a_i \otimes a_i'\) be an element of B satisfying the conditions in item (ii) of Proposition 5. Let L be an extension of the field K; we must prove that the ring \(\mathrm{A}_{(\mathrm{L})}\) is semisimple or, equivalently, that every \(\mathrm{A}_{(\mathrm{L})}\) -module is semisimple (VIII, p.~138, Proposition 4). Let M be an \(\mathrm{A}_{(\mathrm{L})}\) -module, and let N be a submodule of M; we view M as a left (A,L)-bimodule and N as a sub-bimodule (III, §4, No.~3, p.~466). Since L is a field, there exists an L-linear projector \(u\) in M with image N. Since the homotheties \(a_{\mathrm{M}}\) associated with the elements \(a\) of A are L-linear, there exists a unique group homomorphism from \(\mathrm{A} \otimes_{\mathrm{K}} \mathrm{A}^{\circ}\) to \(\operatorname{End}_{\mathrm{L}}(\mathrm{M})\) that sends an element \(a \otimes a'\) to the L-linear mapping \(x \mapsto au(a'x)\) . We denote the image of \(e\) by this homomorphism by \(v\) ; it follows from Remark 2 that \(v\) is an \(\mathrm{A}_{(\mathrm{L})}\) -linear projector with image N. The kernel of \(v\) is an \(\mathrm{A}_{(\mathrm{L})}\) -submodule of M, supplementary to N. By Corollary 2 of VIII, p.~56, the \(\mathrm{A}_{(\mathrm{L})}\) -module M is semisimple.

Remark 3. --- We know (VIII, p.~234, Corollary 1) that the tensor product of two absolutely semisimple algebras over a commutative field is absolutely semisimple. Consequently, if the algebra A is absolutely semisimple, then so is the algebra \(B = A \otimes_{K} A^{\circ}\) .

